\chapter{Lesson 1: Shopping and Money}

CHINA 102R picks up where 101R ended. We assume the 101R vocabulary (about 500 words),
the basic word order (subject -- time -- place -- verb -- object), \zh{的}, measure
words, the degree complement and verb + \zh{了}.

\section{Vocabulary}

\begin{center}
\begin{tabular}{llll}
\toprule
Characters & Pinyin & Part of speech & Meaning \\
\midrule
\zh{买} & mǎi & verb & to buy \\
\zh{卖} & mài & verb & to sell \\
\zh{东西} & dōngxi & noun & things, stuff \\
\zh{商店} & shāngdiàn & noun & store \\
\zh{超市} & chāoshì & noun & supermarket \\
\zh{售货员} & shòuhuòyuán & noun & salesperson \\
\zh{钱} & qián & noun & money \\
\zh{块} & kuài & measure word & dollar, yuan (spoken) \\
\zh{毛} & máo & measure word & 0.1 yuan (spoken) \\
\zh{分} & fēn & measure word & 0.01 yuan \\
\zh{元} & yuán & measure word & yuan (written) \\
\zh{多少钱} & duōshao qián & & how much (money) \\
\zh{一共} & yígòng & adverb & altogether \\
\zh{便宜} & piányi & adjective & cheap \\
\zh{贵} & guì & adjective & expensive \\
\zh{件} & jiàn & measure word & for shirts, jackets, matters \\
\zh{条} & tiáo & measure word & for long, thin things (pants, skirts, fish) \\
\zh{双} & shuāng & measure word & pair \\
\zh{付钱} & fù qián & verb-object & to pay \\
\zh{找钱} & zhǎo qián & verb-object & to give change \\
\zh{现金} & xiànjīn & noun & cash \\
\zh{信用卡} & xìnyòngkǎ & noun & credit card \\
\zh{刷卡} & shuā kǎ & verb-object & to pay by card (``swipe'') \\
\zh{扫码} & sǎo mǎ & verb-object & to scan a QR code \\
\zh{打折} & dǎ zhé & verb-object & to give a discount \\
\zh{试} & shì & verb & to try (on) \\
\zh{还是} & háishi & conjunction & or (in questions) \\
\zh{或者} & huòzhě & conjunction & or (in statements) \\
\zh{需要} & xūyào & verb & to need \\
\bottomrule
\end{tabular}
\end{center}

\section{Chinese money}

The currency of the PRC is the \zh{人民币} \emph{Rénmínbì} (RMB, ``people's money''),
whose unit is the \zh{元} \emph{yuán}. In speech everyone says \zh{块}; prices in shops
are written with \zh{元} or the symbol \textyen.

\begin{center}
\begin{tikzpicture}[b/.style={draw, thick, rounded corners, minimum width=2.6cm, minimum height=1.1cm, align=center}]
  \node[b, fill=green!15] (k) at (0,0) {\zh{一块}\\ (\zh{一元})};
  \node[b, fill=blue!12] (m) at (4.5,0) {\zh{十毛}\\ (\zh{十角})};
  \node[b, fill=orange!15] (f) at (9,0) {\zh{一百分}};
  \node at (2.25,0) {$=$};
  \node at (6.75,0) {$=$};
  \node[font=\small, below=0.2cm of k] {spoken (written)};
  \node[font=\small, below=0.2cm of m, align=center] {\zh{毛} spoken, \zh{角} jiǎo written};
  \node[font=\small, below=0.2cm of f] {\zh{分} (rarely used now)};
\end{tikzpicture}
\end{center}

\subsection{Reading prices}

\begin{center}
\begin{tabular}{lll}
\toprule
Price & Full form & Everyday form \\
\midrule
\textyen 3.50 & \zh{三块五毛} & \zh{三块五} \\
\textyen 12.00 & \zh{十二块} & \zh{十二块(钱)} \\
\textyen 2.05 & \zh{两块零五分} & \zh{两块零五} \\
\textyen 0.80 & \zh{八毛} & \zh{八毛(钱)} \\
\textyen 108 & \zh{一百零八块} & \zh{一百零八} \\
\textyen 250 & \zh{两百五十块} & \zh{两百五} \\
\bottomrule
\end{tabular}
\end{center}

\begin{important}[Rules for prices]
\begin{itemize}
  \item The last unit can be dropped: \zh{三块五} = 3.50, \zh{一百二} = 120.
  \item ``2'' before \zh{块}, \zh{毛}, \zh{百}, \zh{千} is usually \zh{两}: \zh{两块},
        \zh{两毛}, \zh{两百}.
  \item An internal gap is filled by \zh{零}: \zh{一百零五} = 105 (not 150).
\end{itemize}
\end{important}

\begin{aside}
By 2018--2019 cash had nearly disappeared in Chinese cities: people pay with WeChat Pay
(\zh{微信支付}) or Alipay (\zh{支付宝}) by scanning a QR code (\zh{扫码}), even at street
stalls. Visitors without a Chinese bank account had trouble paying for anything. ``250''
(\zh{二百五} èrbǎiwǔ) is slang for an idiot, so avoid prices ending in 250 as gifts!
\end{aside}

\section{Grammar}

\subsection{Asking prices}

\begin{itemize}
  \item \zh{这件衬衫多少钱？}---\zh{一百二十块。} How much is this shirt?
  \item \zh{苹果怎么卖？}---\zh{五块钱一斤。} How are the apples sold?---5 yuan per
        \emph{jin} (half a kilogram).
  \item \zh{一共多少钱？}---\zh{一共八十五块五。} How much altogether?
\end{itemize}

\subsection{\zh{的} as a nominalizer}

Adjective / noun / verb phrase + \zh{的} can stand for ``the one that \ldots'' when the
noun is understood.
\begin{itemize}
  \item \zh{我要大的，不要小的。} I want the big one, not the small one.
  \item \zh{这件太贵了，有便宜的吗？} This one is too expensive; are there cheaper ones?
  \item \zh{卖东西的} = a person who sells things, a vendor.
\end{itemize}

\subsection{\zh{还是} vs \zh{或者}}

Both mean ``or''. \zh{还是} is used in \emph{questions} that ask the listener to choose;
\zh{或者} is used in \emph{statements} that list possibilities.

\begin{itemize}
  \item \zh{你要黄的还是红的？} Do you want the yellow one or the red one?
  \item \zh{黄的或者红的都可以。} Either yellow or red is fine.
  \item \zh{我们星期六或者星期天去买东西吧。} Let's go shopping Saturday or Sunday.
\end{itemize}

\begin{important}[Exception]
After \zh{不知道} ``don't know'' or \zh{不一定}, \zh{还是} is used even in a statement,
because there is an embedded question: \zh{我不知道他要红的还是黄的。}
\end{important}

\subsection{Discounts with \zh{打折}}

\zh{打 X 折} means you pay X tenths of the price, so the number is the percentage you
\emph{pay}, not the one you save.
\begin{itemize}
  \item \zh{打八折} = 80\% of the price = 20\% off.
  \item \zh{打七五折} = 75\% of the price = 25\% off.
  \item \zh{打对折} = half price.
\end{itemize}

\begin{example}[Computing a sale price]
A jacket costs \textyen 400 and the sign says \zh{打六折}. You pay
$400 \times 0.6 = 240$ yuan, i.e.\ \zh{两百四十块}.
\end{example}

\subsection{Pairs and measure words for shopping}

\begin{center}
\begin{tabular}{lll}
\toprule
MW & For & Examples \\
\midrule
\zh{件} jiàn & tops, coats, matters & \zh{一件衬衫, 一件大衣, 一件事} \\
\zh{条} tiáo & long, thin things & \zh{一条裤子, 一条裙子, 一条鱼} \\
\zh{双} shuāng & pairs & \zh{一双鞋, 一双筷子} \\
\zh{斤} jīn & half a kilogram & \zh{两斤苹果} \\
\zh{瓶} píng & bottles & \zh{三瓶水} \\
\bottomrule
\end{tabular}
\end{center}

\section{At the store}

\begin{example}[Buying a shirt]\leavevmode
\begin{itemize}
  \item \textbf{Clerk:} \zh{您想买什么？}
  \item \textbf{Customer:} \zh{我想买一件衬衫。}
  \item \textbf{Clerk:} \zh{您要什么颜色的？白的还是蓝的？}
  \item \textbf{Customer:} \zh{白的。中号的。可以试试吗？}
  \item \textbf{Clerk:} \zh{当然可以。}
  \item \textbf{Customer:} \zh{这件不大也不小，正好。多少钱？}
  \item \textbf{Clerk:} \zh{一百五。今天打八折，一百二。}
  \item \textbf{Customer:} \zh{可以刷卡吗？}
  \item \textbf{Clerk:} \zh{可以。扫码也行。}
\end{itemize}
(\zh{中号} = medium size; \zh{正好} zhènghǎo = just right; \zh{当然} dāngrán = of course.)
\end{example}

\section*{Review}
\begin{itemize}
  \item \zh{块/毛/分} in speech, \zh{元/角/分} in writing; the last unit can be dropped.
  \item \zh{两} before \zh{块, 毛, 百, 千}; \zh{零} fills gaps.
  \item \zh{多少钱？}, \zh{怎么卖？}, \zh{一共}.
  \item Adj/N/VP + \zh{的} = ``the one that \ldots''.
  \item \zh{还是} in questions, \zh{或者} in statements.
  \item \zh{打八折} = pay 80\%.
  \item Measure words \zh{件, 条, 双, 斤}.
\end{itemize}
